\chapter{Twists Of Supersymmetric Theories}
\label{ch:tqft-twists-of-supersymmetric-theories}

\ConventionsEuclidean

A topological twist is a change of the rotation action on fields and
operators using an \(R\)-symmetry of a supersymmetric QFT.  The goal is to
produce a scalar odd symmetry whose cohomology defines metric-independent
observables under stated Ward-identity hypotheses.

\section{Twisting Datum}

Let a Euclidean supersymmetric QFT in dimension \(d\) have rotation group
\(\operatorname{Spin}(d)\) and compact \(R\)-symmetry group \(G_R\).  A
twisting datum is a continuous homomorphism
\[
  \rho:\operatorname{Spin}(d)\longrightarrow G_R .
\]
The twisted rotation group is the diagonal subgroup
\[
  \operatorname{Spin}(d)_{\rho}
  \subset
  \operatorname{Spin}(d)\times G_R,
  \qquad
  g\longmapsto (g,\rho(g)).
\]
Every field representation is restricted to this diagonal subgroup.  A
supercharge \(Q\) is a candidate cohomological differential when it is a
scalar under \(\operatorname{Spin}(d)_\rho\) and obeys
\[
  Q^2=\delta_{\mathrm{gauge}}(\Lambda)+\delta_{\mathrm{flavor}}(\alpha)
\]
on gauge-invariant, flavor-invariant observables.  If translations appear on
the right side, the resulting theory is cohomological only after the
translation is trivial in the chosen observable sector.

\section{Cohomological Observables}

Let \(\mathcal A\) be the graded algebra of gauge-invariant local or extended
observables in the twisted theory.  The odd derivation
\[
  d_Q\mathcal O=[Q,\mathcal O\}_{\mathrm{gr}}
\]
defines the cohomology
\[
  H_Q(\mathcal A)=\ker d_Q/\operatorname{im}d_Q .
\]
Correlation functions of \(Q\)-closed observables descend to \(H_Q(\mathcal A)\)
provided the finite-regulator functional integral or Hilbert-space trace has
no boundary contribution under the \(Q\)-variation.

\begin{proposition}[Metric independence from a cohomological stress tensor]
\label{prop:twist-metric-independence}
Let \(Z_g(\mathcal O_1,\ldots,\mathcal O_n)\) be a regulated correlation
functional on a compact manifold \(M\) for \(Q\)-closed insertions.  Assume
that the regulator preserves \(Q\), that no boundary term arises in the
space of fields, and that the metric variation of the action is
\[
  \frac{\delta S}{\delta g^{\mu\nu}(x)}
  =
  \{Q,G_{\mu\nu}(x)\}_{\mathrm{gr}}
\]
up to contact terms whose integrals vanish against the declared insertions.
Then \(Z_g(\mathcal O_1,\ldots,\mathcal O_n)\) is locally constant as a
function of the metric.
\end{proposition}

\begin{proof}
For a variation \(\delta g^{\mu\nu}\),
\[
  \delta_g Z_g
  =
  -\int_M d^dx\,\delta g^{\mu\nu}(x)
  \left\langle
  \{Q,G_{\mu\nu}(x)\}_{\mathrm{gr}}
  \prod_i\mathcal O_i
  \right\rangle .
\]
Since \(d_Q\mathcal O_i=0\), the integrand is the \(Q\)-variation of the
product \(G_{\mu\nu}(x)\prod_i\mathcal O_i\), up to the contact terms included
in the hypothesis.  The \(Q\)-Ward identity for the regulator sets the
expectation value of this variation to zero.  Hence every first metric
variation vanishes.
\end{proof}

\section{Donaldson Twist}

In four Euclidean dimensions,
\[
  \operatorname{Spin}(4)=SU(2)_+\times SU(2)_- .
\]
The \(\mathcal N=2\) vector multiplet has \(R\)-symmetry
\[
  G_R=SU(2)_R\times U(1)_r .
\]
The Donaldson twist uses the diagonal subgroup
\[
  SU(2)'_+=\operatorname{diag}(SU(2)_+\times SU(2)_R),
  \qquad
  \operatorname{Spin}(4)'=SU(2)'_+\times SU(2)_- .
\]
Under this twisted rotation group the supercharges decompose into a scalar,
a one-form, and a self-dual two-form.  The scalar supercharge is the
cohomological differential.  The twisted vector multiplet fields may be
organized as
\[
  A_\mu,\qquad
  \phi,\bar\phi,\qquad
  \psi_\mu,\qquad
  \chi_{\mu\nu}^+,\qquad
  \eta ,
\]
with auxiliary fields added so that the \(Q\)-transformations close off shell
up to gauge transformations.

\section{Two-Dimensional A And B Twists}

For a two-dimensional theory with \(\mathcal N=(2,2)\) supersymmetry, the
rotation group is \(U(1)_{\mathrm{rot}}\) and the \(R\)-symmetry contains
vector and axial \(U(1)\) factors when they are nonanomalous.  The A twist
uses \(U(1)_V\); the B twist uses \(U(1)_A\).  The A-model cohomological
observables depend on symplectic data and have instanton sectors given by
holomorphic maps.  The B-model observables depend on complex-structure data
and are represented perturbatively by polyvector-field cohomology.  These
statements require the anomaly and compactness assumptions developed in
Chapter~\ref{ch:tqft-topological-sigma-models}.

\section{Regulator And Contact-Term Ledger}

Twisting changes the representation of fields under rotations; it does not
by itself construct a continuum QFT.  A complete cohomological QFT datum must
include:
\[
\begin{array}{ll}
  \text{field space or algebra} & \mathcal F,\\
  \text{odd differential} & Q,\\
  \text{closure relation} & Q^2,\\
  \text{action or state} & S \text{ or } \omega,\\
  \text{measure or trace} & D\Phi \text{ at finite regulator},\\
  \text{contact terms} & \text{operator-collision prescriptions}.
\end{array}
\]
Only after these entries are fixed can metric independence, localization, and
descent equations be treated as mathematical consequences.

\begin{openproblem}[Twisted continuum construction]
\label{op:twisted-continuum-construction}
Construct four-dimensional twisted supersymmetric gauge theories as
continuum QFTs with a \(Q\)-preserving regulator, off-shell closure, contact
term control, and a derivation of the cohomological stress-tensor relation.
\end{openproblem}
